A translationally invariant PEPS lives on a torus: a finite square lattice whose
opposite boundaries are identified, so that the two coordinates are read modulo
the periods $n$ and $m$.
\begin{center}
    % Formula/source: Papers/1804.04964/paper_normal.tex, lines 104--112,
    % defines a translationally invariant PEPS by repeating one tensor on a
    % regular lattice and contracting according to the lattice geometry;
    % lines 1451--1471 state the normal-PEPS theorem on an n-by-m region.
    % Ink-to-index map: the nine upper legs are the free physical indices of
    % the displayed 3-by-3 window.  Every horizontal and vertical line is a
    % virtual bond; the marked A is a representative of the tensor repeated
    % at every site.
    % Contraction set: twelve nearest-neighbour bonds inside the window are
    % summed, and six further bonds identify the three west/east endpoint
    % pairs and the three north/south endpoint pairs.  No virtual index stays
    % open, while all nine physical indices remain free.
    % Boundary signature: virtual=0 | physical=9.
    % Source verdict: the source has no separate torus-geometry picture.  This
    % is the house-style realization of the chapter's cyclic-coordinate
    % definition; unlike the former open-boundary schematic, it displays the
    % periodic contraction in both coordinate directions.
    \begin{tenkzlattice}[
        rows=3, cols=3, sheets={ket}, outer legs=top, frame=oblique,
        west=trace, east=trace, north=trace, south=trace]
        \tnsite[role=marked, label=$A$]{(1,2)}
    \end{tenkzlattice}
\end{center}
Translations act on the vertices and edges. Transporting a tensor along a graph
isomorphism $\phi$ relabels its coefficients by
$\Coeff_{A^\phi}(\sigma)=\Coeff_A(\sigma\circ\phi)$,
so two tensors generating the same state still do after transport. This is the
setting in which the normal Fundamental Theorem specializes to a single
site-independent tensor and produces the global phase of
\cite[Theorem~3]{Molnar2018NormalPEPS}.

\begin{definition}[Torus geometry and translations]%
    \label{def:peps_torus_geometry_translation}
    \lean{TNLean.PEPS.TorusVertex}
    \lean{TNLean.PEPS.torusGraph}
    \lean{TNLean.PEPS.torusGraph_adj_right}
    \lean{TNLean.PEPS.torusGraph_adj_up}
    \lean{TNLean.PEPS.Edge.map}
    \lean{TNLean.PEPS.Edge.equiv}
    \lean{TNLean.PEPS.IncidentEdge.equiv}
    \lean{TNLean.PEPS.translateEquiv}
    \lean{TNLean.PEPS.translate}
    \lean{TNLean.PEPS.translateEdge}
    \lean{TNLean.PEPS.translateEdge_isHorizontal}
    \lean{TNLean.PEPS.translateEdge_isVertical}
    \lean{TNLean.PEPS.translateIncidentEdge}
    \lean{TNLean.PEPS.torusEdge_horizontal_or_vertical}
    \leanok
    The discrete $n\times m$ torus has vertex set $\Z/n\Z\times\Z/m\Z$.
    Its graph joins $(x,y)$ to $(x+1,y)$ and to $(x,y+1)$, with the
    coordinates read cyclically; every edge is of one of these two forms,
    so every edge is horizontal or vertical. Translation by $(a,b)$ is the
    graph automorphism $\tau_{a,b}(x,y)=(x+a,y+b)$.
    The induced action on edges sends an edge with unordered endpoint pair
    $\{u,v\}$ to the edge with unordered endpoint pair
    $\{\tau_{a,b}(u),\tau_{a,b}(v)\}$; in particular it preserves the
    horizontal and vertical orientation classes. For a general graph
    isomorphism $\phi:G\simeq G'$, the same endpoint action gives a
    bijection on edges and a bijection from the edges incident to $v$ to
    the edges incident to $\phi(v)$.
\end{definition}

\begin{definition}[Transport of a PEPS tensor]%
    \label{def:peps_tensor_transport}
    \lean{TNLean.PEPS.Tensor.transport}
    \lean{TNLean.PEPS.Tensor.transport_bondDim}
    \leanok
    \uses{def:peps_tensor, def:peps_torus_geometry_translation}
    Let $\phi:G\simeq G'$ be a graph isomorphism. The transported tensor
    $A^\phi$ on $G'$ has bond dimension
    $D^\phi_{e'}=D_{\phi^{-1}(e')}$ at every edge $e'$ of $G'$. Its local
    tensor at $w\in G'$ is the local tensor of $A$ at $\phi^{-1}(w)$, with
    virtual indices moved through the incident-edge bijection induced by
    $\phi$.
\end{definition}

\begin{theorem}[State coefficient under transport]%
    \label{thm:peps_stateCoeff_transport}
    \lean{TNLean.PEPS.stateCoeff_transport}
    \leanok
    \uses{def:peps_stateCoeff, def:peps_tensor_transport}
    For every graph isomorphism $\phi:G\simeq G'$ and every physical
    configuration $\sigma$ on the vertices of $G'$,
    $\Coeff_{A^\phi}(\sigma)=\Coeff_A(\sigma\circ\phi)$.
\end{theorem}
\begin{proof}\leanok
    The edge bijection induced by $\phi$ identifies virtual configurations
    of $A^\phi$ with virtual configurations of $A$. After this change of
    summation variables, the local factor at $w\in G'$ becomes the local
    factor of $A$ at $\phi^{-1}(w)$. Reindexing the finite product over
    vertices by $\phi$ gives the coefficient equality.
\end{proof}

\begin{theorem}[Transport preserves equality of PEPS states]%
    \label{thm:peps_sameState_transport}
    \lean{TNLean.PEPS.SameState.transport}
    \leanok
    \uses{def:peps_same_state, thm:peps_stateCoeff_transport}
    If two PEPS tensors on $G$ have equal state coefficients for every
    physical configuration, then their transports along any graph
    isomorphism $\phi:G\simeq G'$ have equal state coefficients for every
    physical configuration on $G'$.
\end{theorem}
\begin{proof}\leanok
    By Theorem~\ref{thm:peps_stateCoeff_transport},
    \begin{align}
        \Coeff_{A^\phi}(\sigma)
        &=\Coeff_A(\sigma\circ\phi)
          =\Coeff_B(\sigma\circ\phi)
          =\Coeff_{B^\phi}(\sigma),
        \notag
    \end{align}
    where the middle equality is the assumed equality of state coefficients.
\end{proof}

\begin{definition}[Translation-invariant torus tensor]%
    \label{def:peps_torus_translation_invariant}
    \lean{TNLean.PEPS.IsTorusTranslationInvariant}
    \leanok
    \uses{def:peps_tensor_transport, def:peps_torus_geometry_translation}
    A PEPS tensor $A$ on the discrete torus is translation invariant if
    $A^{\tau_{a,b}}=A$ for every translation $\tau_{a,b}$. This formalizes
    the setting of \cite[Theorem~3]{Molnar2018NormalPEPS}, where one tensor
    is repeated at every site of a periodic lattice.
\end{definition}

\begin{theorem}[State coefficient of a translation-invariant torus tensor]%
    \label{thm:peps_stateCoeff_translationInvariant}
    \lean{TNLean.PEPS.stateCoeff_translationInvariant}
    \leanok
    \uses{def:peps_torus_translation_invariant,
        thm:peps_stateCoeff_transport}
    If $A$ is translation invariant on the discrete torus, then
    $\Coeff_A(\sigma\circ\tau_{a,b})=\Coeff_A(\sigma)$ for every translation
    $\tau_{a,b}$ and every physical configuration $\sigma$.
\end{theorem}
\begin{proof}\leanok
    Theorem~\ref{thm:peps_stateCoeff_transport} gives
    \begin{align}
        \Coeff_A(\sigma\circ\tau_{a,b})
        &=\Coeff_{A^{\tau_{a,b}}}(\sigma)
          =\Coeff_A(\sigma),
        \notag
    \end{align}
    where the second equality substitutes $A^{\tau_{a,b}}=A$ from translation
    invariance.
\end{proof}

\begin{definition}[Orientation-uniform bond dimensions on the torus]%
    \label{def:peps_torus_uniform_bond_dim}
    \lean{TNLean.PEPS.torusRightEdge}
    \lean{TNLean.PEPS.torusUpEdge}
    \lean{TNLean.PEPS.TorusUniformBondDim}
    \leanok
    \uses{def:peps_torus_geometry_translation}
    A bond-dimension function on the torus is orientation-uniform with
    horizontal dimension $D_h$ and vertical dimension $D_v$ if
    $D_{e_h}=D_h$ for every horizontal edge $e_h$ and $D_{e_v}=D_v$ for
    every vertical edge $e_v$. The reference horizontal edge is the right
    edge at the origin, and the reference vertical edge is the up edge at
    the origin.
\end{definition}

\begin{theorem}[Translation invariance gives orientation-uniform bond dimensions]%
    \label{thm:peps_torusUniformBondDim_of_translationInvariant}
    \lean{TNLean.PEPS.translateEdge_eq_map}
    \lean{TNLean.PEPS.bondDim_translateEdge_of_translationInvariant}
    \lean{TNLean.PEPS.isHorizontalTorusEdge_eq_rightEdge}
    \lean{TNLean.PEPS.isVerticalTorusEdge_eq_upEdge}
    \lean{TNLean.PEPS.translateEdge_torusRightEdge}
    \lean{TNLean.PEPS.translateEdge_torusUpEdge}
    \lean{TNLean.PEPS.bondDim_torusRightEdge_const}
    \lean{TNLean.PEPS.bondDim_torusUpEdge_const}
    \lean{TNLean.PEPS.torusUniformBondDim_of_translationInvariant}
    \leanok
    \uses{def:peps_torus_translation_invariant,
        def:peps_torus_uniform_bond_dim}
    If $A$ is translation invariant on the discrete torus, then its bond
    dimensions are orientation-uniform. Explicitly, with
    $D_h=D_{e_h}$ at the reference right edge $e_h$ and
    $D_v=D_{e_v}$ at the reference up edge $e_v$, one has
    $D_{f_h}=D_h$ for every horizontal edge $f_h$ and $D_{f_v}=D_v$ for
    every vertical edge $f_v$.
\end{theorem}
\begin{proof}\leanok
    Translation invariance gives $A^{\tau_{a,b}}=A$, hence
    $D_A(\tau_{a,b}(e))=D_A(e)$ for every edge $e$. Every horizontal edge
    $f_h$ has the form $f_h=\tau_{a,b}(e_h)$ for a translate of the reference
    right edge, so
    \begin{align}
        D_{f_h}
        &=D_A(\tau_{a,b}(e_h))
          =D_A(e_h)
          =D_h.
        \notag
    \end{align}
    The same argument with the reference up edge gives $D_{f_v}=D_v$.
\end{proof}

\paragraph{Failure of the row-and-column reduction.}

\begin{definition}[Per-edge product matrix on two vertical bonds]%
    \label{def:peps_row_cut_per_edge_product}
    \lean{TNLean.PEPS.RowColumnReductionObstruction.IsPerEdgeProduct}
    \leanok
    A matrix $M\in \MN{4}$ is a per-edge product on two collected vertical
    bonds if, under the identification
    $(r,s)\longmapsto 2r+s$ between
    $\{0,1\}\times\{0,1\}$ and $\{0,1,2,3\}$, there are matrices
    $Y_0,Y_1\in \MN{2}$ such that, for all
    $r,s,r',s'\in\{0,1\}$,
    \begin{align}
        M_{(r,s),(r',s')}
        &=(Y_0)_{r,r'}(Y_1)_{s,s'}.
        \notag
    \end{align}
\end{definition}

\begin{definition}[Per-edge factorization of the row-cut gauge]%
    \label{def:peps_row_cut_gauge_factorizes}
    \lean{TNLean.PEPS.RowColumnReductionObstruction.RowCutGaugeFactorizes}
    \leanok
    \uses{def:mps_tensor, def:l_blk_injective, def:same_mpv,
        def:peps_row_cut_per_edge_product}
    The row-cut factorization assertion is the following statement. Let $A$
    and $B$ be MPS tensors with physical alphabet $\{0,\ldots,15\}$ and
    bond space $\C^4$. If $A$ and $B$ are $1$-block injective and generate
    the same MPV family, then there exist $Z\in \GL_4(\C)$ and
    $\lambda\in\C$ such that $Z$ is a per-edge product matrix and
    $B^i=\lambda Z^{-1}A^iZ$ for all physical indices
    $i\in\{0,\ldots,15\}$. This is the additional product conclusion needed
    by a row-wise one-dimensional reduction before it could imply the per-edge
    conclusion of \cite[Theorem~3]{Molnar2018NormalPEPS}.
\end{definition}

\begin{theorem}[A row-cut gauge need not factor per edge]%
    \label{thm:peps_row_cut_gauge_factorizes_false}
    \lean{TNLean.PEPS.RowColumnReductionObstruction.rowCutGaugeFactorizes_false}
    \leanok
    \uses{def:peps_row_cut_gauge_factorizes}
    The row-cut factorization assertion is false. Explicitly, let
    $E_{ab}\in\MN{4}$ be the matrix units and set
    \begin{align}
        A^{(a,b)}
        &=E_{ab}, \notag\\
        G
        &=
        \begin{pmatrix}
            1&0&1&0\\
            0&1&0&0\\
            0&0&1&0\\
            0&0&0&1
        \end{pmatrix}, \notag\\
        B^{(a,b)}
        &=G^{-1}E_{ab}G.
        \notag
    \end{align}
    Then $A$ and $B$ are $1$-block injective and generate the same MPV
    family, but there are no $Z\in\GL_4(\C)$ and $\lambda\in\C$ with
    $Z$ a per-edge product matrix and
    $B^{(a,b)}=\lambda Z^{-1}E_{ab}Z$ for all
    $a,b\in\{0,\ldots,3\}$.
\end{theorem}

\begin{proof}
    \leanok
    The matrix units span $\MN{4}$, hence $A$ is $1$-block injective; conjugation
    by $G$ preserves this property for $B$. For every word
    $w=((a_1,b_1),\ldots,(a_N,b_N))$, one has
    $B^w=G^{-1}A^wG$ and $\tr(B^w)=\tr(A^w)$, so the two MPV families are
    equal.

    Suppose that $Z$ and $\lambda$ satisfy
    $B^{(a,b)}=\lambda Z^{-1}E_{ab}Z$ for every $a,b\in\{0,\ldots,3\}$.
    With $W=ZG^{-1}$, multiplication by $Z$ on the left and by $G^{-1}$
    on the right gives $WE_{ab}=\lambda E_{ab}W$ for every matrix unit
    $E_{ab}$. Taking the $(x,y)$ entry gives
    $W_{x,a}\delta_{b,y}=\lambda\delta_{x,a}W_{b,y}$ for all
    $x,y,a,b\in\{0,\ldots,3\}$. With $y=b$, this gives
    $x\ne a\Rightarrow W_{x,a}=0$ and
    $W_{a,a}=\lambda W_{b,b}$.

    Since $Z$ and $G$ are invertible, so is $W$. Thus the diagonal entries
    cannot all vanish; the case $a=b$ gives $\lambda=1$, and then
    $W_{a,a}=W_{b,b}$ for all $a,b$. Hence $W=c\Id_4$ and $\lambda=1$ for
    some $c\ne0$, so $Z=cG$.

    For a matrix $M\in\MN{4}$, write $M_{rr'}\in\MN{2}$ for the block
    whose $(s,s')$ entry is $M_{(r,s),(r',s')}$. If $Z$ were a per-edge
    product, then $Z_{rr'}=(Y_0)_{r,r'}Y_1$ for some
    $Y_0,Y_1\in\MN{2}$. In particular,
    $(Z_{00})_{s,s'}(Z_{01})_{1,1}
    =(Z_{01})_{s,s'}(Z_{00})_{1,1}$.
    Since $Z=cG$ and $c\ne0$, scaling cancels from this block
    proportionality condition. The corresponding blocks of $G$ are
    \begin{align}
        G_{00}
        &=
        \begin{pmatrix}
            1&0\\
            0&1
        \end{pmatrix}, \notag\\
        G_{01}
        &=
        \begin{pmatrix}
            1&0\\
            0&0
        \end{pmatrix}.
        \notag
    \end{align}
    Substitution into
    $(G_{00})_{00}(G_{01})_{11}
    =(G_{01})_{00}(G_{00})_{11}$ gives $0=1$. Hence no per-edge product
    conjugator can exist.
\end{proof}


\paragraph{Staircase-window geometry for a reference edge.}

\begin{lemma}[Left endpoint of the staircase reference edge]%
    \label{lem:peps_reference_edge_left_endpoint_staircase_window}
    \lean{TNLean.PEPS.referenceEdge_leftEndpoint_mem_staircaseWindow}
    \leanok
    Let $L,K>0$, let $s=(a,b)$ be a staircase corner on the torus, and assume
    $a+2L\le n$ and $2K\le m$. Let $e$ be the horizontal reference edge
    from $(a+L-1,b+K-1)$ to $(a+L,b+K-1)$, and let $W_j$ be the $j$th
    staircase window based at $s$. Then the left endpoint
    $u=(a+L-1,b+K-1)$ lies in $W_j$ if and only if $1\le j$.
\end{lemma}
\begin{proof}\leanok
    The endpoint formula gives the two coordinate distances from the corner:
    the column distance is $L-1$ and the row distance is $K-1$. Substituting
    these distances into the staircase-window membership criterion gives exactly
    $1\le j$.
\end{proof}

\begin{lemma}[Right endpoint of the staircase reference edge]%
    \label{lem:peps_reference_edge_right_endpoint_staircase_window}
    \lean{TNLean.PEPS.referenceEdge_rightEndpoint_mem_staircaseWindow}
    \leanok
    In the same notation, the right endpoint $w=(a+L,b+K-1)$ lies in $W_j$
    if and only if $j<L$.
\end{lemma}
\begin{proof}\leanok
    The right endpoint has column distance $L$ and row distance $K-1$ from
    $s$. The row condition is automatic, while the column band contains this
    distance precisely for $j<L$.
\end{proof}

\begin{lemma}[Boundary windows for the staircase reference edge]%
    \label{lem:peps_reference_edge_boundary_windows}
    \lean{TNLean.PEPS.isRegionBoundaryEdge_staircaseWindow_referenceEdge}
    \leanok
    \uses{lem:peps_reference_edge_left_endpoint_staircase_window,
        lem:peps_reference_edge_right_endpoint_staircase_window}
    If $j=0$ or $L\le j$, then the reference edge $e$ is a boundary edge
    of $W_j$: exactly one of its two endpoints lies in the window.
\end{lemma}
\begin{proof}\leanok
    For $j=0$, Lemmas~\ref{lem:peps_reference_edge_left_endpoint_staircase_window}
    and~\ref{lem:peps_reference_edge_right_endpoint_staircase_window} show that
    only $w$ lies in $W_j$. For $L\le j$, the same two endpoint tests show
    that only $u$ lies in $W_j$.
\end{proof}

\begin{lemma}[Interior windows for the staircase reference edge]%
    \label{lem:peps_reference_edge_interior_windows}
    \lean{TNLean.PEPS.referenceEdge_endpoints_mem_staircaseWindow_of_interior}
    \leanok
    \uses{lem:peps_reference_edge_left_endpoint_staircase_window,
        lem:peps_reference_edge_right_endpoint_staircase_window}
    If $1\le j<L$, then both endpoints of $e$ lie in $W_j$.
\end{lemma}
\begin{proof}\leanok
    The left-endpoint criterion gives $u\in W_j$ from $1\le j$, and the
    right-endpoint criterion gives $w\in W_j$ from $j<L$.
\end{proof}

\begin{lemma}[Interior windows are not boundary windows]%
    \label{lem:peps_reference_edge_not_boundary_interior_windows}
    \lean{TNLean.PEPS.not_isRegionBoundaryEdge_staircaseWindow_referenceEdge_of_interior}
    \leanok
    \uses{lem:peps_reference_edge_interior_windows}
    If $1\le j<L$, then $e$ is not a boundary edge of $W_j$.
\end{lemma}
\begin{proof}\leanok
    By Lemma~\ref{lem:peps_reference_edge_interior_windows}, both endpoints of
    $e$ lie in $W_j$. The definition of a boundary edge requires exactly one
    endpoint in the region, so $e$ is not a boundary edge of $W_j$.
\end{proof}
